\section{Human Robot Interaction Challenge}
\label{test:human-robot-interaction-challenge}

\subsection*{Description}
The robot assists people during a social event taking place throughout the
\Arena{}. Several people are present, but only some of them require assistance.
The robot must explore the house, recognize natural requests for attention,
approach only people requesting assistance, and understand and fulfil their
requests.

\textbf{Main goal:}
The robot autonomously discovers who requires assistance and completes requests
through natural, attentive human--robot interaction. It must not rely on knowing
in advance who needs help, where that person is located, or what they will ask.

\textbf{Interaction types:}
The task focuses on natural human--robot interaction rather than general-purpose
command understanding and execution. Interactions may require the robot to:
\begin{enumerate}[nosep]
	\item provide information during a conversation;
	\item request an object from the barman and deliver it to the requester.
\end{enumerate}

\subsection*{Focus}
\emph{\SysI{}}, \emph{\HRI{}}, \emph{person detection},
\emph{gesture and sound-source recognition}, \emph{social attention},
\emph{knowledge questions}, and \emph{navigation}.

\subsection*{Setup}
\begin{itemize}
	\item \textbf{Location:}
	\begin{itemize}
		\item The test takes place throughout the domestic environment during a
		social event or party.
		\item The robot starts outside the \Arena{} and enters when the entrance
		door is opened.
	\end{itemize}

	\item \textbf{People:}
	\begin{itemize}
		\item Seven designated requesters are distributed naturally across several
		rooms and may be sitting, standing, talking, or performing activities
		appropriate to the event.
		\item Additional people may be present in the house as bystanders. The
		number, roles, and locations of all people are not announced to the team.
		\item Both designated requesters and bystanders may move naturally around
		the house, as people would during a normal social event.
		\item A person requiring assistance calls or waves to attract the robot's
		attention only while the robot is in the same room as that person.
		\item People not requesting assistance behave as bystanders and must not be
		approached as requesters.
	\end{itemize}

	\item \textbf{Requests:}
	\begin{itemize}
		\item The content and location of requests are not announced to the team
		before the test.
		\item Requests are part of the interaction scenarios described below. The
		specific interaction variants are defined by the TC and OC.
		\item Objects used by an interaction scenario follow the applicable object
		and arena rules.
	\end{itemize}

	\item \textbf{Barman table:}
	\begin{itemize}
		\item A barman table containing all objects that may be requested is located
		inside the \Arena{}.
		\item A barman stands behind the table. The robot does not search for or pick
		up requested objects itself; the barman places the requested object on the
		robot.
	\end{itemize}
\end{itemize}

\subsection*{Interaction types}
Each request uses one of the following two interaction types and is presented
through one of the interaction scenarios defined below.
\begin{enumerate}[nosep]
	\item \textbf{Provide information:} The robot answers a question during the
	conversation. Questions are limited to objects and object classes, rooms and
	furniture inside the \Arena{}, the date and time, the competition, or the
	robot and its team, including the team's history and past events in which the
	team participated.
	\item \textbf{Deliver an object:} The robot goes to the barman table, informs
	the barman which object was requested, and asks the barman to place it on the
	robot. The barman may not be paying attention. The robot must confirm, in any
	way it considers suitable, that the object is in its possession before
	returning to the requester and delivering the object.
\end{enumerate}

\subsection*{Procedure}
\begin{enumerate}[nosep]
	\item \textbf{Test start:} The entrance door is opened and the robot enters
	the \Arena{}.
	\item \textbf{Discover people requesting assistance:} The robot explores the
	house and attends to the people present. It must determine whether a person is
	requesting assistance from natural social cues. The robot must only approach a
	person as a requester after that person has requested its attention.
	\item \textbf{Understand the request:} The robot approaches the requester and
	interacts with them to determine what assistance is required.
	\item \textbf{Complete the interaction:} The robot responds according to the
	assigned interaction scenario, maintaining attention and conversational
	context as required by that scenario.
	\item \textbf{Continue assisting:} After completing an interaction, the robot
	continues exploring and looking for other people requesting assistance until
	the test ends.
\end{enumerate}

\subsection*{Interaction scenarios}
The seven requesters take part in six different interaction scenarios. The
two-person drink request involves two requesters; each other scenario involves
one requester.
\begin{enumerate}[nosep]
	\item \textbf{Follow while conversing:} A requester asks the robot to follow
	them. While the robot is following, the requester starts and maintains a
	conversation. The robot must follow safely while remaining attentive and
	responding to the requester.
	\item \textbf{Full-duplex conversation:} The requester interacts without
	waiting for rigid speaking turns. The person may respond while the robot is
	speaking or naturally interrupt it. The robot must recognize the interruption,
	yield or adapt its response, and continue the conversation without losing the
	relevant context.
	\item \textbf{Correcting a request:} The requester gives the robot a command.
	As the robot starts moving to execute it, the requester calls the robot, asks
	it to stop, and corrects or changes part of the command. The robot must stop
	safely, attend to the requester, understand the correction, and continue using
	the updated command instead of the original one.
	\item \textbf{Two simultaneous drink requests:} Two requesters ask the robot
	for different drinks at approximately the same time. The robot must manage the
	overlapping requests, determine what each person asked for, associate each
	drink with the correct requester, communicate both orders to the barman, and
	deliver each drink to the correct person.
	\item \textbf{Local-language interaction:} The requester speaks a language
	used in the region hosting the competition. The robot must understand the
	request and answer in the same language. The local language used for this
	interaction is selected by the OC.
	\item \textbf{Learning a new drink:} A requester teaches the robot about a
	drink that is not part of the announced known or common object lists. The
	person provides information that uniquely distinguishes the drink. The
	requester then asks for that drink. The robot must use the information learned
	during the interaction to communicate the correct order to the barman and
	deliver the drink provided by the barman.
\end{enumerate}

\subsection*{Additional rules and remarks}
\begin{enumerate}[nosep]
	\item \textbf{No advance request assignment:} Teams are not told which people
	will request assistance, their positions, or their requests. The robot must use
	a generic strategy to discover and handle interactions.
	\item \textbf{Request for attention:} A natural request for attention may be
	verbal or non-verbal. Calling the robot and waving are canonical examples. A
	requester calls or waves only while the robot is in the same room; therefore,
	the robot must explore the house rather than wait for a global call. Detailed
	cues and interaction variants will be specified by the TC and OC.
	\item \textbf{Do not disturb bystanders:} The robot may perceive and observe
	all people while exploring, but it must not initiate a task-oriented
	conversation with a person who has not requested assistance.
	\item \textbf{Dynamic gaze:} During a conversation, the robot must direct its
	gaze toward the person with whom it is interacting. If that person moves, the
	robot must dynamically adjust its gaze to remain attentive. Brief gaze shifts
	that are natural or necessary for safe operation are allowed.
\end{enumerate}

\subsection*{Instructions}

\subsubsection*{To Referee}
The referees need to:
\begin{itemize}
	\item ensure that only designated requesters use a request-for-attention cue;
	\item ensure that bystanders behave naturally without misleading the robot;
	\item give each requester the assigned request privately; and
	\item brief the requesters on their assigned interaction scenario and ensure
	they perform it consistently; and
	\item ensure that the barman follows the prescribed interaction and places
	the requested object on the robot without requiring the robot to pick it up.
\end{itemize}

\subsubsection*{To OC}
Before the test, the OC needs to:
\begin{itemize}
	\item recruit seven requesters and enough additional volunteers to act as
	bystanders, as well as one barman, and distribute them naturally across the
	house;
	\item provide the barman table and place all objects that may be requested on
	it;
	\item prepare the information, objects, and other materials required by the
	six interaction scenarios, including the local language and
	the unannounced drink used for interactive learning; and
	\item ensure that the selected requests can be executed and scored reliably
	without revealing requester identities or locations to the teams.
\end{itemize}

\subsection*{Score sheet}
\begin{scorelist}[startbutton=false,timelimit=7]
	\scoreheading{Request for attention}
	\scoreitem[7]{50}{Approach a person calling the robot}
	\scoreitem[7]{300}{Correctly solve an interaction scenario\newline
	(might be a good idea to specify for each scenario because of different difficulty levels)}
		\scorepen{50}{Fail to dynamically look at the person speaking}
		\scoremod{50}{Correctly answer an information question}
		\scoremod{50}{Correctly inform the barman of the requested item}
		\scoremod{100}{Deliver a requested object to the correct person}
		\scorepen{100}{Leave the barman without the requested item}
		\scorepen{100}{Deliver a requested object to the incorrect person}

	\scoreheading{Penalties}
	\penaltyitem{50}{Approach a person who did not call the robot (per person)}
\end{scorelist}


% Local Variables:
% TeX-master: "Rulebook"
% End:


% Local Variables:
% TeX-master: "Rulebook"
% End:
